%% =============================================================================
%% chapters/appendix-mllp-services.tex
%% Appendix B: Implementation Guide for MLLP Ingress and Egress Services
%% =============================================================================

\chapter{Implementation Guide for MLLP Ingress and Egress Services}
\label{app:mllp_services}

\section{Overview \& Engineering Objectives}
\label{sec:mllp_guide_overview}

This implementation guide provides an exhaustive, production-grade technical manual detailing the architectural patterns, component blueprints, configuration parameters, state persistence workflows, and testing harnesses required to implement:
\begin{enumerate}
    \item \textbf{A New MLLP Ingress Service}: Ingesting inbound HL7 v2 messages over Minimal Lower Layer Protocol (MLLP) TCP sockets, validating syntax, constructing canonical FHIR R5 \texttt{Communication} and \texttt{Task} resources, persisting records into the \texttt{Mneme} Infinispan cache grid and \texttt{Mnemosyne} PostgreSQL database, publishing asynchronous \texttt{ErgonEvent} notifications to the \texttt{Petasos} (ActiveMQ Artemis) message bus, and generating synchronous HL7 acknowledgements (\texttt{AA}, \texttt{AE}, \texttt{AR}).
    \item \textbf{A New MLLP Egress Service}: Consuming outbound clinical dispatch tasks from dedicated \texttt{Petasos} destination queues, transforming canonical models into HL7 v2 payloads, dispatching over Apache Camel MLLP client routes with robust socket timeouts, deeply inspecting downstream HL7 \texttt{ACK}/\texttt{NACK} responses, managing full FHIR \texttt{Communication}, \texttt{Task}, and \texttt{Provenance} lifecycle state transitions in \texttt{Mneme}/\texttt{Mnemosyne}, and isolating destination failure domains through multi-instance container topologies.
\end{enumerate}

% =============================================================================
% PART I: IMPLEMENTING AN MLLP INGRESS SERVICE
% =============================================================================
\section{Ingress Architectural Blueprint \& Framing Protocols}
\label{sec:ingress_blueprint}

The MLLP Ingress Gateway (\texttt{pylai-mllp-in}) acts as the secure, high-throughput perimeter gateway of Harmonia, terminating inbound TCP connections from external Patient Administration Systems (PAS), Laboratory Information Systems (LIS), Radiology Information Systems (RIS), and Electronic Medical Records (EMR).

\begin{figure}[htbp]
    \centering
    \begin{adjustbox}{max width=\textwidth}
        %% =============================================================================
%% diagrams/fig-mllp-ingress-process.tex
%% ArchiMate 3.2 MLLP Ingress Processing Architecture & Component Flow View
%% =============================================================================
\begin{tikzpicture}[node distance=1.0cm and 0.8cm, auto]

    % --- Column 1: External Source & MLLP Network Layer ---
    \node[archi-business-actor] (ext_source) at (0, 7.5) {%
        \archibadge{Business Actor}\archiname{Clinical Source}\\\archidesc{PAS / LIS / RIS / EMR}%
    };

    \node[archi-app-interface] (mllp_endpoint) at (0, 4.5) {%
        \archibadge{App Interface}\archiname{Camel MLLP Endpoint}\\\archidesc{mllp://0.0.0.0:2101--2104}%
    };

    \node[archi-app-service] (sync_ack_service) at (0, 1.5) {%
        \archibadge{App Service}\archiname{Synchronous ACK Engine}\\\archidesc{AA / AE / AR Generation}%
    };

    % --- Column 2: Ingress Routing & Processing Layer ---
    \node[archi-app-component] (route_builder) at (4.5, 7.5) {%
        \archibadge{App Component}\archiname{MLLP RouteBuilder}\\\archidesc{Incoming*RouteBuilder}%
    };

    \node[archi-app-component] (proc_wrapper) at (4.5, 4.5) {%
        \archibadge{App Component}\archiname{Message Processor}\\\archidesc{Incoming*ProcessorWrapper}%
    };

    \node[archi-app-process] (hl7_extractor) at (4.5, 1.5) {%
        \archibadge{App Process}\archiname{HL7 Data Extractor}\\\archidesc{HAPI Terser / Validation}%
    };

    % --- Column 3: FHIR Transformation & Persistence Layer ---
    \node[archi-app-process] (fhir_builders) at (9.0, 7.5) {%
        \archibadge{App Process}\archiname{FHIR Resource Builders}\\\archidesc{Communication \& Task}%
    };

    \node[archi-app-component] (cache_service) at (9.0, 4.5) {%
        \archibadge{App Component}\archiname{Mneme \& Mnemosyne}\\\archidesc{HotRod Cache \& JPA Store}%
    };

    \node[archi-data-object] (canonical_resources) at (9.0, 1.5) {%
        \archibadge{Data Object}\archiname{FHIR R5 Resources}\\\archidesc{Communication \& Task}%
    };

    % --- Column 4: Asynchronous Messaging & Event Dispatch Layer ---
    \node[archi-app-component] (event_producer) at (13.5, 7.5) {%
        \archibadge{App Component}\archiname{TaskEventProducer}\\\archidesc{TaskEventProducerService}%
    };

    \node[archi-app-interface] (petasos_queue) at (13.5, 4.5) {%
        \archibadge{App Interface}\archiname{Petasos Artemis Queue}\\\archidesc{petasos.queue.inbound.*}%
    };

    \node[archi-app-component] (ponos_engine) at (13.5, 1.5) {%
        \archibadge{App Component}\archiname{Ponos Workflow Engine}\\\archidesc{Ergon Pipeline Dispatch}%
    };

    % --- Flow & Structural Relationships ---
    % Ingress Inbound TCP Socket Flow
    \draw[archi-flow] (ext_source) -- node[left, font=\tiny] {HL7 MLLP Frame} (mllp_endpoint);
    \draw[archi-serving] (mllp_endpoint) -- (route_builder);
    \draw[archi-triggering] (route_builder) -- (proc_wrapper);
    \draw[archi-assignment] (proc_wrapper) -- (hl7_extractor);

    % Extraction to FHIR Resource Construction
    \draw[archi-triggering] (hl7_extractor) -- (fhir_builders);
    \draw[archi-access] (fhir_builders) -- (canonical_resources);

    % Persistence Flows
    \draw[archi-flow] (fhir_builders) -- (cache_service);
    \draw[archi-access] (cache_service) -- (canonical_resources);

    % Petasos Messaging Dispatch
    \draw[archi-triggering] (cache_service) -- (event_producer);
    \draw[archi-flow] (event_producer) -- node[right, font=\tiny] {ErgonEvent / TaskEvent} (petasos_queue);
    \draw[archi-flow] (petasos_queue) -- (ponos_engine);

    % Synchronous ACK Return Loop
    \draw[archi-triggering] (cache_service) to[bend left=20] node[above, font=\tiny] {Success / Error Signal} (sync_ack_service);
    \draw[archi-flow] (sync_ack_service) -- node[left, font=\tiny] {HL7 ACK (AA/AE/AR)} (ext_source);

\end{tikzpicture}

    \end{adjustbox}
    \caption{ArchiMate 3.2 Component Flow for the MLLP Ingress Service Processing Pipeline}
    \label{fig:mllp_ingress_process}
\end{figure}

\subsection{MLLP Transport Framing}

MLLP provides minimal framing over stream-oriented TCP sockets to delineate discrete message boundaries. All inbound MLLP ingress routes adhere to the standard framing specification:
\begin{equation}
\langle\text{Inbound Frame}\rangle ::= \texttt{0x0B}\ (\text{Start-of-Block / \texttt{<VT>}}) + \langle\text{HL7 v2.x Payload}\rangle + \texttt{0x1C}\ (\text{End-of-Block / \texttt{<FS>}}) + \texttt{0x0D}\ (\text{Carriage Return / \texttt{<CR>}})
\end{equation}

Inbound frames are received over persistent or transient TCP sockets. Camel's Netty-based MLLP component buffers incoming bytes until the End-of-Block trailer (\texttt{0x1C 0x0D}) is detected, strips framing bytes, and passes the raw HL7 string to the downstream pipeline.

\subsection{The 8-Phase Ingress Processing Lifecycle}

Every inbound message traverses an immutable 8-phase execution pipeline:
\begin{enumerate}
    \item \textbf{Socket Connection \& TCP Frame Capture}: The client establishes a TCP connection to the configured port (e.g., 2101 for ADT, 2102 for Lab ORU, 2103 for Rad ORU, 2104 for ORM). Camel captures the framed byte stream.
    \item \textbf{Frame Stripping \& Charset Normalization}: Framing bytes (\texttt{0x0B}, \texttt{0x1C}, \texttt{0x0D}) are stripped; bytes are decoded into UTF-8 or ISO-8859-1 strings based on \texttt{MSH-18} character set headers.
    \item \textbf{Message Type \& Trigger Event Validation}: The MSH segment is inspected to identify message structure (\texttt{MSH-9.1}), trigger event (\texttt{MSH-9.2}), and message control ID (\texttt{MSH-10}).
    \item \textbf{HL7 v2 Data Extraction}: Specialized extractors leverage HAPI Terser to navigate segment hierarchies and extract demographic, clinical, and encounter fields.
    \item \textbf{Canonical FHIR R5 Construction}: Canonical resource builders instantiate:
    \begin{itemize}
        \item A FHIR R5 \texttt{Communication} resource preserving the complete raw HL7 message in \texttt{Communication.payload.contentAttachment}.
        \item A synthetic FHIR R5 \texttt{Task} resource (representing the executable \texttt{Pragma}) encapsulating workflow priority, routing metadata, and correlation IDs.
    \end{itemize}
    \item \textbf{Cache \& Database Persistence}: The \texttt{Communication} and \texttt{Task} resources are persisted into the \texttt{Mneme} Infinispan data grid via \texttt{TaskCacheService} and write-behind synced to \texttt{Mnemosyne} JPA stores.
    \item \textbf{Petasos Event Dispatch}: A lightweight \texttt{ErgonEvent} notification referencing the persistent \texttt{Task.id} and \texttt{Communication.id} is published to the target ActiveMQ Artemis inbound queue (e.g., \texttt{petasos.queue.inbound.adt}).
    \item \textbf{Synchronous HL7 Acknowledgement Generation}: A standards-compliant HL7 \texttt{ACK} containing matching \texttt{MSA-2} and appropriate \texttt{MSA-1} code (\texttt{AA}, \texttt{AE}, \texttt{AR}) is framed in MLLP and flushed over the TCP socket before closing the connection.
\end{enumerate}

\section{Ingress Component Implementation}
\label{sec:ingress_components}

\subsection{Camel RouteBuilder Configuration}

Ingress routes are defined using Apache Camel's Java DSL, configured as CDI \texttt{@Dependent} beans. The route registers an \texttt{mllp://} endpoint, sets connection options, and wires incoming exchanges to dedicated processor wrappers.

\begin{lstlisting}[language=Java, caption={IncomingAdtMessageMllpRouteBuilder Implementation Pattern}]
@Dependent
public class IncomingAdtMessageMllpRouteBuilder extends RouteBuilder {

    private static final Logger log = LoggerFactory.getLogger(IncomingAdtMessageMllpRouteBuilder.class);

    @Inject
    private MllpConfig mllpConfig;

    @Inject
    private IncomingAdtMessageProcessorWrapper incomingAdtMessageProcessorWrapper;

    @Override
    public void configure() throws Exception {
        if (mllpConfig == null) {
            mllpConfig = new MllpConfig();
        }

        String host = mllpConfig.getHost();
        int port = mllpConfig.getAdtPort(); // Port 2101
        boolean autoAck = mllpConfig.isAutoAck(); // false: explicit ACK control
        int bindTimeout = mllpConfig.getBindTimeout();
        int readTimeout = mllpConfig.getReadTimeout();

        // MLLP Ingress Listener Route
        from(String.format("mllp://%s:%d?autoAck=%b&bindTimeout=%d&readTimeout=%d",
                host, port, autoAck, bindTimeout, readTimeout))
            .routeId("hl7-mllp-adt-receiver")
            .log(LoggingLevel.INFO, log, "MLLP Inbound ADT message received on port " + port)
            .process(incomingAdtMessageProcessorWrapper)
            .log(LoggingLevel.INFO, log, "MLLP Inbound ADT processing completed; ACK ready");

        // Internal Direct Route for In-Process Triggers
        from("direct:adt-events")
            .routeId("direct-adt-processor")
            .process(incomingAdtMessageProcessorWrapper);
    }
}
\end{lstlisting}

\subsection{Processor Wrapper \& Context Management}

The \texttt{IncomingAdtMessageProcessorWrapper} acts as the boundary adapter between Apache Camel's exchange model and Harmonia's HL7 domain processor.

\begin{lstlisting}[language=Java, caption={IncomingAdtMessageProcessorWrapper Blueprint}]
@ApplicationScoped
public class IncomingAdtMessageProcessorWrapper implements Processor {

    private static final Logger log = LoggerFactory.getLogger(IncomingAdtMessageProcessorWrapper.class);

    @Inject
    private IncomingAdtMessageProcessor incomingAdtMessageProcessor;

    @Override
    public void process(Exchange exchange) throws Exception {
        String hl7Message = exchange.getIn().getBody(String.class);
        log.debug("Processing inbound HL7 payload of length: {}", hl7Message.length());

        // Delegate to domain processor
        AdtProcessingResult result = incomingAdtMessageProcessor.process(hl7Message);

        // Populate exchange headers and set synchronous ACK body
        exchange.getMessage().setHeader("HIE_MESSAGE_CONTROL_ID", result.getMessageControlId());
        exchange.getMessage().setHeader("HIE_ACK_CODE", result.getAckCode());
        exchange.getMessage().setBody(result.getAckMessage());
    }
}
\end{lstlisting}

\subsection{HL7 v2 Data Extraction with HAPI Terser}

To ensure resilience against optional or missing segments across varied HL7 versions (2.3, 2.4, 2.5), extractors employ the HAPI HL7 Terser with defensive navigation wrappers:

\begin{lstlisting}[language=Java, caption={AdtMessageExtractor Implementation Blueprint}]
public class AdtMessageExtractor {

    private final Terser terser;

    public AdtMessageExtractor(Message message) {
        this.terser = new Terser(message);
    }

    public String extractMessageControlId() {
        return getSafeValue("/.MSH-10");
    }

    public String extractTriggerEvent() {
        return getSafeValue("/.MSH-9-2");
    }

    public String extractPatientId() {
        String pid3 = getSafeValue("/.PID-3-1");
        return (pid3 != null && !pid3.isEmpty()) ? pid3 : getSafeValue("/.PID-2-1");
    }

    public String extractFamilyName() {
        return getSafeValue("/.PID-5-1");
    }

    public String extractGivenName() {
        return getSafeValue("/.PID-5-2");
    }

    public String extractDateOfBirth() {
        return getSafeValue("/.PID-7-1");
    }

    public String extractAdministrativeGender() {
        return getSafeValue("/.PID-8");
    }

    public String extractVisitNumber() {
        return getSafeValue("/.PV1-19-1");
    }

    public String extractPatientClass() {
        return getSafeValue("/.PV1-2");
    }

    public String extractAssignedLocation() {
        String poc = getSafeValue("/.PV1-3-1"); // Point of care
        String room = getSafeValue("/.PV1-3-2");
        String bed = getSafeValue("/.PV1-3-3");
        return String.format("%s/%s/%s", 
            poc != null ? poc : "", 
            room != null ? room : "", 
            bed != null ? bed : "");
    }

    private String getSafeValue(String terserPath) {
        try {
            String value = terser.get(terserPath);
            return (value != null && !value.trim().isEmpty()) ? value.trim() : null;
        } catch (HL7Exception e) {
            return null; // Defensive fallback on missing segment/field
        }
    }
}
\end{lstlisting}

\section{Canonical FHIR Model Mapping \& Resource Builders}
\label{sec:ingress_fhir_mapping}

Harmonia normalizes all inbound HL7 messages into dual canonical FHIR R5 entities:
\begin{enumerate}
    \item \textbf{FHIR R5 \texttt{Communication}}: Holds message provenance, sender/recipient metadata, and the raw immutable HL7 v2 payload.
    \item \textbf{FHIR R5 \texttt{Task}}: Models the executable units of work (\texttt{Pragma}) consumed by the \texttt{Ponos} workflow engine.
\end{enumerate}

\subsection{FHIR Communication Resource Construction}

\begin{lstlisting}[language=Java, caption={AdtCommunicationResourceBuilder Blueprint}]
public class AdtCommunicationResourceBuilder {

    public Communication buildCommunication(AdtMessageExtractor extractor, String rawHl7) {
        Communication comm = new Communication();
        comm.setId(UUID.randomUUID().toString());
        comm.setStatus(Communication.CommunicationStatus.COMPLETED);

        // Identifiers
        Identifier msh10Id = new Identifier();
        msh10Id.setSystem("http://harmonia.hie/identifiers/msh-10");
        msh10Id.setValue(extractor.extractMessageControlId());
        comm.addIdentifier(msh10Id);

        // Category
        CodeableConcept category = new CodeableConcept();
        category.addCoding().setSystem("http://terminology.hl7.org/CodeSystem/communication-category")
                            .setCode("notification")
                            .setDisplay("HL7 Inbound Notification");
        comm.addCategory(category);

        // Sender & Recipient
        comm.setSender(new Reference().setDisplay(extractor.getSendingFacility()));
        comm.addRecipient(new Reference().setDisplay("Harmonia-HIE-Ingress"));

        // Payload Attachment (Raw HL7 Message)
        Communication.CommunicationPayloadComponent payload = new Communication.CommunicationPayloadComponent();
        Attachment attachment = new Attachment();
        attachment.setContentType("text/plain");
        attachment.setData(rawHl7.getBytes(StandardCharsets.UTF_8));
        attachment.setCreation(new Date());
        payload.setContent(attachment);
        comm.addPayload(payload);

        comm.setSent(new Date());
        return comm;
    }
}
\end{lstlisting}

\subsection{FHIR Task (Pragma) Resource Construction}

The synthetic \texttt{Task} represents the executable pipeline instance. Table~\ref{tab:hl7_task_mapping} outlines the canonical field mappings.

\begin{table}[htbp]
\centering
\small
\begin{tabularx}{\textwidth}{l p{3.5cm} l p{4.5cm}}
\toprule
\textbf{HL7 v2 Source Field} & \textbf{FHIR R5 Task Property} & \textbf{Datatype} & \textbf{Transformation Logic} \\
\midrule
\texttt{MSH-10} & \texttt{Task.identifier[msh10]} & \texttt{Identifier} & Copied verbatim for correlation \\
\texttt{MSH-9-2} (Trigger) & \texttt{Task.businessStatus} & \texttt{CodeableConcept} & Mapped to clinical event type (e.g. \texttt{A01} $\to$ \texttt{ADMISSION}) \\
\texttt{PV1-2} (Patient Class) & \texttt{Task.priority} & \texttt{TaskPriority} & \texttt{E} $\to$ \texttt{stat}, \texttt{I} $\to$ \texttt{urgent}, \texttt{O} $\to$ \texttt{routine} \\
\texttt{PID-3-1} (MRN) & \texttt{Task.for} & \texttt{Reference} & Reference to \texttt{Patient/MRN-\{id\}} \\
\texttt{PV1-19-1} (Visit Num) & \texttt{Task.encounter} & \texttt{Reference} & Reference to \texttt{Encounter/VIS-\{num\}} \\
Auto-generated UUID & \texttt{Task.id} & \texttt{IdType} & Unique Pragma identifier \\
Literal \texttt{"requested"} & \texttt{Task.status} & \texttt{TaskStatus} & Initial state for Ponos engine \\
Literal \texttt{"order"} & \texttt{Task.intent} & \texttt{TaskIntent} & Required FHIR workflow intent \\
\bottomrule
\end{tabularx}
\caption{HL7 v2 to FHIR R5 Task Canonical Mapping Specifications}
\label{tab:hl7_task_mapping}
\end{table}

\section{Persistence, Event Dispatch \& Synchronous Acknowledgement}
\label{sec:ingress_persistence_ack}

\subsection{Persistence Integration with Mneme \& Mnemosyne}

The ingress processor writes canonical entities into the \texttt{Mneme} Infinispan data grid via \texttt{TaskCacheService} and \texttt{CommunicationService}:

\begin{lstlisting}[language=Java, caption={Ingress Persistence \& Event Dispatch Implementation}]
public AdtProcessingResult process(String rawHl7) {
    try {
        Message hl7Msg = hapiParser.parse(rawHl7);
        AdtMessageExtractor extractor = new AdtMessageExtractor(hl7Msg);

        // 1. Build Canonical Resources
        Communication comm = commBuilder.buildCommunication(extractor, rawHl7);
        Task task = taskBuilder.buildTask(extractor, comm.getId());

        // 2. Persist to Mneme Cache Grid / Mnemosyne Store
        communicationService.save(comm);
        taskService.save(task);

        // 3. Publish Asynchronous Notification to Petasos Message Bus
        ErgonEvent event = new ErgonEvent();
        event.setEventId(UUID.randomUUID().toString());
        event.setTaskId(task.getId());
        event.setCommunicationId(comm.getId());
        event.setMessageType("ADT");
        event.setTriggerEvent(extractor.extractTriggerEvent());
        event.setTimestamp(System.currentTimeMillis());

        taskEventProducerService.publishInboundEvent("petasos.queue.inbound.adt", event);

        // 4. Return Synchronous Application Accept (AA)
        String ackHl7 = hl7AckGenerator.generateAck(hl7Msg, "AA", "Message Accepted");
        return new AdtProcessingResult(extractor.extractMessageControlId(), "AA", ackHl7);

    } catch (HL7Exception e) {
        log.error("HL7 parse error: {}", e.getMessage());
        String nackHl7 = hl7AckGenerator.generateNack(rawHl7, "AR", "Structural Parse Failure: " + e.getMessage());
        return new AdtProcessingResult(null, "AR", nackHl7);
    } catch (Exception e) {
        log.error("Internal processing error: {}", e.getMessage());
        String nackHl7 = hl7AckGenerator.generateNack(rawHl7, "AE", "Internal Storage Failure: " + e.getMessage());
        return new AdtProcessingResult(null, "AE", nackHl7);
    }
}
\end{lstlisting}

\subsection{Synchronous Acknowledgement Mapping Rules}

The ingress service guarantees deterministic acknowledgement status codes as summarized in Table~\ref{tab:ack_rules}.

\begin{table}[htbp]
\centering
\small
\begin{tabularx}{\textwidth}{l l p{6.0cm} p{3.0cm}}
\toprule
\textbf{ACK Code} & \textbf{Status Name} & \textbf{Triggering Condition} & \textbf{Client Action} \\
\midrule
\texttt{AA} & Application Accept & Syntactically valid HL7; successfully written to Mneme cache and Petasos queue. & Message considered durably accepted; do not retry. \\
\texttt{AE} & Application Error & Transient infrastructure failure (Infinispan timeout, database disconnect, Artemis broker unreachable). & Re-queue and retry with exponential backoff. \\
\texttt{AR} & Application Reject & Malformed HL7 frame, missing required MSH segments, unsupported major version. & Route to operator dead-letter queue; do not retry. \\
\bottomrule
\end{tabularx}
\caption{Synchronous HL7 Acknowledgement Response Contracts}
\label{tab:ack_rules}
\end{table}

\subsection{REST Management \& CDI Deployment Configuration}

Every ingress gateway exposes optional REST endpoints for health checks, administrative ping, and direct HTTP ingestion:

\begin{lstlisting}[language=Java, caption={JAX-RS Direct Ingestion Endpoint (AdtResource)}]
@Path("/mllp/adt")
@ApplicationScoped
public class AdtResource {

    @Inject
    private IncomingAdtMessageProcessor processor;

    @POST
    @Consumes(MediaType.TEXT_PLAIN)
    @Produces(MediaType.TEXT_PLAIN)
    public Response ingestAdtMessage(String rawHl7) {
        AdtProcessingResult result = processor.process(rawHl7);
        if ("AA".equals(result.getAckCode())) {
            return Response.ok(result.getAckMessage()).build();
        } else if ("AE".equals(result.getAckCode())) {
            return Response.status(Response.Status.SERVICE_UNAVAILABLE).entity(result.getAckMessage()).build();
        } else {
            return Response.status(Response.Status.BAD_REQUEST).entity(result.getAckMessage()).build();
        }
    }
}
\end{lstlisting}

% =============================================================================
% PART II: IMPLEMENTING AN MLLP EGRESS SERVICE
% =============================================================================
\section{Multi-Instance Egress Architecture \& Queue Partitioning}
\label{sec:egress_architecture}

In contrast to ingress gateways that bind to inbound listening ports, the MLLP Egress Gateway (\texttt{pylai-mllp-out}) functions as a dynamic TCP client, connecting to downstream clinical endpoints (such as hospital EMR, LMS, and RIS-PACS systems).

\begin{figure}[htbp]
    \centering
    \begin{adjustbox}{max width=\textwidth}
        %% =============================================================================
%% diagrams/fig-mllp-egress-process.tex
%% ArchiMate 3.2 MLLP Egress Multi-Instance Architecture & Lifecycle View
%% =============================================================================
\begin{tikzpicture}[node distance=1.0cm and 0.8cm, auto]

    % --- Column 1: Petasos Artemis Messaging Layer ---
    \node[archi-app-interface] (petasos_dest_queue) at (0, 7.5) {%
        \archibadge{App Interface}\archiname{Petasos Destination Queue}\\\archidesc{petasos.queue.egress.*}%
    };

    \node[archi-app-component] (queue_consumer) at (0, 4.5) {%
        \archibadge{App Component}\archiname{Task Queue Consumer}\\\archidesc{OutboundTaskQueueConsumer}%
    };

    \node[archi-app-interface] (petasos_comp_queue) at (0, 1.5) {%
        \archibadge{App Interface}\archiname{Completion Queue}\\\archidesc{petasos.queue.egress.completed}%
    };

    % --- Column 2: Egress Processing & Camel Route Layer ---
    \node[archi-app-component] (outbound_proc) at (4.5, 7.5) {%
        \archibadge{App Component}\archiname{Outbound Processor}\\\archidesc{OutboundMllpProcessor}%
    };

    \node[archi-app-component] (camel_producer) at (4.5, 4.5) {%
        \archibadge{App Component}\archiname{Camel Dynamic Producer}\\\archidesc{OutboundMllpRouteBuilder (toD)}%
    };

    \node[archi-app-component] (ack_proc) at (4.5, 1.5) {%
        \archibadge{App Component}\archiname{HL7 ACK Processor}\\\archidesc{Hl7AckProcessor (MSA/ERR)}%
    };

    % --- Column 3: Downstream Network & Target Clinical System ---
    \node[archi-app-interface] (target_mllp_socket) at (9.0, 7.5) {%
        \archibadge{App Interface}\archiname{Remote TCP Socket}\\\archidesc{mllp://host:2201--2205}%
    };

    \node[archi-business-actor] (target_system) at (9.0, 4.5) {%
        \archibadge{Business Actor}\archiname{Downstream Receiver}\\\archidesc{EMR / LMS / RIS-PACS}%
    };

    \node[archi-data-object] (ack_payload) at (9.0, 1.5) {%
        \archibadge{Data Object}\archiname{HL7 ACK / NACK Frame}\\\archidesc{AA, AE, AR, CA, CE, CR}%
    };

    % --- Column 4: Lifecycle, Persistence & Provenance Layer ---
    \node[archi-app-component] (lifecycle_mgr) at (13.5, 7.5) {%
        \archibadge{App Component}\archiname{Lifecycle Manager}\\\archidesc{OutboundStateLifecycleManager}%
    };

    \node[archi-app-process] (provenance_builder) at (13.5, 4.5) {%
        \archibadge{App Process}\archiname{Provenance Builder}\\\archidesc{OutboundProvenanceResourceBuilder}%
    };

    \node[archi-app-component] (mneme_storage) at (13.5, 1.5) {%
        \archibadge{App Component}\archiname{Mneme \& Mnemosyne}\\\archidesc{Task, Comm \& Provenance Store}%
    };

    % --- Flow & Structural Relationships ---
    % JMS Consumption & Request Preparation
    \draw[archi-flow] (petasos_dest_queue) -- node[left, font=\tiny] {JMS TextMessage} (queue_consumer);
    \draw[archi-triggering] (queue_consumer) -- (outbound_proc);
    \draw[archi-triggering] (outbound_proc) -- node[left, font=\tiny] {direct:mllp-outbound-send} (camel_producer);

    % Outbound MLLP Dispatch over TCP
    \draw[archi-flow] (camel_producer) -- (target_mllp_socket);
    \draw[archi-flow] (target_mllp_socket) -- node[right, font=\tiny] {MLLP Frame} (target_system);

    % Return ACK from Target System
    \draw[archi-flow] (target_system) -- (ack_payload);
    \draw[archi-flow] (ack_payload) -- (ack_proc);
    \draw[archi-triggering] (camel_producer) -- (ack_proc);

    % Lifecycle, Provenance & Persistence Updates
    \draw[archi-triggering] (ack_proc) to[bend left=20] node[above, font=\tiny] {OutboundMllpResponse} (lifecycle_mgr);
    \draw[archi-assignment] (lifecycle_mgr) -- (provenance_builder);
    \draw[archi-flow] (provenance_builder) -- (mneme_storage);
    \draw[archi-access] (lifecycle_mgr) -- (mneme_storage);

    % Completion Event Dispatch
    \draw[archi-flow] (lifecycle_mgr) to[bend right=30] node[above, font=\tiny] {Petasos Completion Event} (petasos_comp_queue);

\end{tikzpicture}

    \end{adjustbox}
    \caption{ArchiMate 3.2 Component Flow for the MLLP Egress Multi-Instance Dispatch Pipeline}
    \label{fig:mllp_egress_process}
\end{figure}

\subsection{Head-of-Line Blocking Mitigation \& Destination Isolation}

In healthcare integration environments, downstream clinical systems exhibit disparate network latency, maintenance windows, and failure modes. If a single outbound gateway serviced all destinations from a shared queue:
\begin{enumerate}
    \item A transient network outage or TCP socket stall at one hospital endpoint (e.g., LMS analyzer reboot) would block outbound transmissions to healthy endpoints (e.g., EMR, RIS-PACS).
    \item Retry storms and queue head-of-line blocking would cascade across the entire healthcare exchange.
\end{enumerate}

To eliminate head-of-line blocking, Harmonia enforces \textbf{destination-partitioned multi-instance isolation}:
\begin{itemize}
    \item Each downstream destination is allocated a dedicated, persistent ActiveMQ Artemis queue:
    \begin{equation}
    \text{Queue Pattern} ::= \texttt{petasos.queue.egress.}\langle\text{destinationId}\rangle\texttt{.}\langle\text{messageType}\rangle
    \end{equation}
    Examples: \texttt{petasos.queue.egress.emr.adt}, \texttt{petasos.queue.egress.lms.orm}, \texttt{petasos.queue.egress.rispac.orm}.
    \item Each destination runs as an isolated, containerized \texttt{pylai-mllp-out} microservice instance consuming exclusively from its dedicated queue.
    \item Failures, exponential backoff retries, and socket timeouts in one container never impact message throughput or delivery SLAs of other destinations.
\end{itemize}

\subsection{Outbound Configuration \& Destination Registry}

Configuration is managed via \texttt{MllpOutboundConfig} and \texttt{MllpDestinationRegistry}, populated via microprofile environment variables at container launch.

\begin{table}[htbp]
\centering
\small
\begin{tabularx}{\textwidth}{l p{4.0cm} l p{4.0cm}}
\toprule
\textbf{Environment Variable} & \textbf{Config Property} & \textbf{Default} & \textbf{Description} \\
\midrule
\texttt{MLLP\_OUTBOUND\_DESTINATION\_ID} & \texttt{destinationId} & \texttt{"DEFAULT"} & Unique destination identifier (e.g. \texttt{EMR}, \texttt{LMS}) \\
\texttt{MLLP\_OUTBOUND\_HOST} & \texttt{host} & \texttt{"localhost"} & Target TCP hostname or IP address \\
\texttt{MLLP\_OUTBOUND\_PORT} & \texttt{port} & \texttt{2575} & Target MLLP TCP listener port \\
\texttt{MLLP\_OUTBOUND\_READ\_TIMEOUT\_MS} & \texttt{readTimeoutMs} & \texttt{10000} & Socket read timeout in milliseconds \\
\texttt{MLLP\_OUTBOUND\_CONNECT\_TIMEOUT\_MS} & \texttt{connectTimeoutMs} & \texttt{5000} & TCP handshake timeout in milliseconds \\
\texttt{MLLP\_OUTBOUND\_AUTO\_ACK} & \texttt{autoAck} & \texttt{false} & Requires explicit downstream ACK verification \\
\texttt{MLLP\_OUTBOUND\_QUEUE\_NAME} & \texttt{dedicatedEventQueueName} & \textit{Derived} & Dedicated Artemis subscription queue \\
\texttt{MLLP\_OUTBOUND\_COMPLETION\_QUEUE\_NAME} & \texttt{completionEventQueueName} & \texttt{"petasos.queue.egress.completed"} & Completion telemetry queue \\
\bottomrule
\end{tabularx}
\caption{MLLP Outbound Gateway Environment Configuration Parameters}
\label{tab:outbound_config_params}
\end{table}

\subsection{Outbound Data Transfer Models}

The egress pipeline coordinates transmissions using two standard DTOs:
\begin{itemize}
    \item \textbf{\texttt{OutboundMllpRequest}}: Encapsulates \texttt{taskId}, \texttt{communicationId}, \texttt{destinationId}, \texttt{messageControlId}, \texttt{targetHost}, \texttt{targetPort}, \texttt{rawHl7Payload}, and retry metadata.
    \item \textbf{\texttt{OutboundMllpResponse}}: Encapsulates delivery status (\texttt{isSuccessful}), \texttt{ackCode} (\texttt{AA}, \texttt{AE}, \texttt{AR}, \texttt{CA}, \texttt{CE}, \texttt{CR}), \texttt{rawAckMessage}, \texttt{errorMessage}, \texttt{durationMs}, and \texttt{acknowledgedAt} timestamp.
\end{itemize}

\section{Outbound Queue Consumer \& Camel Dynamic MLLP Routes}
\label{sec:egress_routes}

\subsection{Dedicated Task Queue Consumer}

The \texttt{OutboundTaskQueueConsumer} implements JMS \texttt{MessageListener}, pulling outbound tasks from the dedicated Artemis queue and dispatching them through Apache Camel.

\begin{lstlisting}[language=Java, caption={OutboundTaskQueueConsumer Implementation Pattern}]
@ApplicationScoped
public class OutboundTaskQueueConsumer implements MessageListener {

    private static final Logger log = LoggerFactory.getLogger(OutboundTaskQueueConsumer.class);

    @Inject
    private MllpOutboundConfig outboundConfig;

    @Inject
    private OutboundStateLifecycleManager lifecycleManager;

    private ProducerTemplate producerTemplate;

    @Override
    public void onMessage(jakarta.jms.Message message) {
        try {
            if (message instanceof TextMessage) {
                String jsonPayload = ((TextMessage) message).getText();
                OutboundMllpRequest request = JsonUtils.fromJson(jsonPayload, OutboundMllpRequest.class);

                // Dispatch via Camel ProducerTemplate
                OutboundMllpResponse response = producerTemplate.requestBody(
                    "direct:mllp-outbound-send", request, OutboundMllpResponse.class);

                // Update FHIR Task, Communication and Provenance states
                lifecycleManager.handleOutboundResponse(request, response);
                message.acknowledge(); // Acknowledge JMS message upon lifecycle recording
            }
        } catch (Exception e) {
            log.error("Fatal exception in outbound queue consumer: {}", e.getMessage(), e);
            // Non-acknowledged message triggers Artemis redelivery policy
        }
    }
}
\end{lstlisting}

\subsection{Camel Dynamic Outbound Route Builder}

The \texttt{OutboundMllpRouteBuilder} configures dynamic dispatch to external TCP sockets using Camel's \texttt{toD} dynamic URI expression, setting connection timeouts and auto-acknowledgment rules.

\begin{lstlisting}[language=Java, caption={OutboundMllpRouteBuilder Dynamic Dispatch DSL}]
@ApplicationScoped
public class OutboundMllpRouteBuilder extends RouteBuilder {

    private static final Logger log = LoggerFactory.getLogger(OutboundMllpRouteBuilder.class);

    @Inject
    private OutboundMllpProcessor outboundMllpProcessor;

    @Inject
    private Hl7AckProcessor hl7AckProcessor;

    @Override
    public void configure() throws Exception {

        // Global Exception Handler for Socket / Network Failures
        onException(Exception.class)
            .handled(true)
            .setHeader("HIE_TRANSMISSION_SUCCESS", constant(false))
            .setHeader("HIE_ACK_CODE", constant("AE"))
            .setHeader("HIE_ERROR_MESSAGE", simple("${exception.message}"))
            .process(exchange -> {
                Exception cause = exchange.getProperty(Exchange.EXCEPTION_CAUGHT, Exception.class);
                String dest = exchange.getIn().getHeader("HIE_DESTINATION_ID", String.class);
                String msh10 = exchange.getIn().getHeader("HIE_EXPECTED_CONTROL_ID", String.class);
                long duration = System.currentTimeMillis() - exchange.getProperty("HIE_START_TIME", Long.class);
                
                OutboundMllpResponse failure = OutboundMllpResponse.failure(
                    msh10, "Socket connection error to " + dest + ": " + cause.getMessage(), duration);
                failure.setDestinationId(dest);
                exchange.getMessage().setBody(failure);
            });

        // Outbound Dynamic MLLP Route
        from("direct:mllp-outbound-send")
            .routeId("mllp-outbound-dynamic-sender")
            .process(outboundMllpProcessor) // Validates request and populates Camel headers
            .setProperty("HIE_START_TIME", simple("${date:now:getTime()}"))
            .log(LoggingLevel.INFO, log, "Dispatching MLLP frame to ${header.HIE_DEST_HOST}:${header.HIE_DEST_PORT}")
            
            // Dynamic MLLP Producer
            .toD("mllp://${header.HIE_DEST_HOST}:${header.HIE_DEST_PORT}?"
                + "autoAck=${header.HIE_AUTO_ACK}"
                + "&readTimeout=${header.HIE_READ_TIMEOUT}"
                + "&connectTimeout=${header.HIE_CONNECT_TIMEOUT}"
                + "&charsetName=UTF-8")
                
            .process(hl7AckProcessor) // Deep inspection of received ACK/NACK frame
            .log(LoggingLevel.INFO, log, "Outbound MLLP dispatch complete: ACK=${header.HIE_ACK_CODE}, Success=${header.HIE_TRANSMISSION_SUCCESS}");
    }
}
\end{lstlisting}

\subsection{Outbound MLLP Processor Header Injection}

The \texttt{OutboundMllpProcessor} validates the inbound DTO, resolves target connection parameters from the \texttt{MllpDestinationRegistry}, and initializes exchange headers:

\begin{lstlisting}[language=Java, caption={OutboundMllpProcessor Header Initialization}]
@ApplicationScoped
public class OutboundMllpProcessor implements Processor {

    @Inject
    private MllpOutboundConfig config;

    @Inject
    private MllpDestinationRegistry destinationRegistry;

    @Override
    public void process(Exchange exchange) throws Exception {
        OutboundMllpRequest request = exchange.getIn().getBody(OutboundMllpRequest.class);
        if (request == null || StringUtils.isBlank(request.getRawHl7Payload())) {
            throw new IllegalArgumentException("OutboundMllpRequest payload must not be null or blank");
        }

        String destinationId = StringUtils.isNotBlank(request.getDestinationId()) 
            ? request.getDestinationId() : config.getDestinationId();

        MllpDestinationConfig destConfig = destinationRegistry.getDestination(destinationId)
            .orElse(config.toDestinationConfig());

        // Configure Exchange Headers for Dynamic toD URI
        exchange.getMessage().setHeader("HIE_DESTINATION_ID", destinationId);
        exchange.getMessage().setHeader("HIE_DEST_HOST", destConfig.getHost());
        exchange.getMessage().setHeader("HIE_DEST_PORT", destConfig.getPort());
        exchange.getMessage().setHeader("HIE_READ_TIMEOUT", destConfig.getReadTimeoutMs());
        exchange.getMessage().setHeader("HIE_CONNECT_TIMEOUT", destConfig.getConnectTimeoutMs());
        exchange.getMessage().setHeader("HIE_AUTO_ACK", false);
        exchange.getMessage().setHeader("HIE_EXPECTED_CONTROL_ID", request.getMessageControlId());

        // Set Outbound Body to Raw HL7 Message String
        exchange.getMessage().setBody(request.getRawHl7Payload());
    }
}
\end{lstlisting}

\section{Downstream HL7 ACK/NACK Deep Inspection}
\label{sec:egress_ack_inspection}

A critical requirement in clinical integration is preventing false-positive transmission confirmations. Receiving a TCP \texttt{FIN} or \texttt{ACK} at the transport layer does not imply clinical acceptance by the receiving hospital system. The \texttt{Hl7AckProcessor} performs deep structural inspection on the received HL7 acknowledgement frame.

\subsection{MSA \& ERR Segment Parsing Specification}

The processor evaluates the \texttt{MSA} (Message Acknowledgement) and \texttt{ERR} (Error) segments:

\begin{lstlisting}[language=Java, caption={Hl7AckProcessor MSA Evaluation Logic}]
public OutboundMllpResponse parseAck(String rawAck, String expectedControlId, String destinationId, long durationMs) {
    if (StringUtils.isBlank(rawAck)) {
        return OutboundMllpResponse.failure(expectedControlId, "Empty ACK response received from " + destinationId, durationMs);
    }

    Matcher matcher = MSA_PATTERN.matcher(rawAck);
    if (!matcher.find()) {
        OutboundMllpResponse resp = OutboundMllpResponse.failure(
            expectedControlId, "Malformed ACK: missing MSA segment", durationMs);
        resp.setRawAckMessage(rawAck);
        return resp;
    }

    String ackCode = matcher.group(1).trim().toUpperCase(); // MSA-1: AA, AE, AR, CA, CE, CR
    String msaControlId = matcher.group(2);                  // MSA-2: Referenced MSH-10
    String ackText = matcher.group(3);                      // MSA-3: Text explanation

    // Strict Correlation Validation
    boolean controlIdMatches = true;
    if (StringUtils.isNotBlank(expectedControlId) && StringUtils.isNotBlank(msaControlId)) {
        if (!StringUtils.equals(expectedControlId.trim(), msaControlId.trim())) {
            controlIdMatches = false;
            log.warn("MSA-2 control ID mismatch: expected '{}' but received '{}'", expectedControlId, msaControlId);
        }
    }

    OutboundMllpResponse response = new OutboundMllpResponse();
    response.setMessageControlId(msaControlId != null ? msaControlId : expectedControlId);
    response.setDestinationId(destinationId);
    response.setAckCode(ackCode);
    response.setAckText(ackText);
    response.setRawAckMessage(rawAck);
    response.setDurationMs(durationMs);
    response.setAcknowledgedAt(new Date());

    boolean isAccept = "AA".equals(ackCode) || "CA".equals(ackCode);
    if (isAccept && controlIdMatches) {
        response.setSuccessful(true);
    } else if (isAccept && !controlIdMatches) {
        response.setSuccessful(false);
        response.setErrorMessage("Control ID mismatch: expected " + expectedControlId + " but received " + msaControlId);
    } else {
        response.setSuccessful(false);
        response.setErrorMessage(StringUtils.isNotBlank(ackText) ? ackText : "HL7 Application NACK (" + ackCode + ")");
    }

    return response;
}
\end{lstlisting}

\subsection{Downstream ACK Code Decision Matrix}

Table~\ref{tab:egress_ack_decision} defines the egress gateway behavior across all downstream ACK codes and error conditions.

\begin{table}[htbp]
\centering
\small
\begin{tabularx}{\textwidth}{l l c p{4.5cm} p{2.5cm}}
\toprule
\textbf{Response Code} & \textbf{Semantics} & \textbf{Success?} & \textbf{State Transition} & \textbf{Retry Strategy} \\
\midrule
\texttt{AA} / \texttt{CA} & Application Accept & \textbf{YES} & \texttt{Task.status} $\to$ \texttt{COMPLETED} & None \\
\texttt{AE} / \texttt{CE} & Application Error & \textbf{NO} & \texttt{Task.status} $\to$ \texttt{IN\_PROGRESS} & Exponential Backoff \\
\texttt{AR} / \texttt{CR} & Application Reject & \textbf{NO} & \texttt{Task.status} $\to$ \texttt{FAILED} & Terminal (DLQ) \\
\texttt{TIMEOUT} & Socket Read Timeout & \textbf{NO} & \texttt{Task.status} $\to$ \texttt{IN\_PROGRESS} & Immediate Reconnect \\
\texttt{CONN\_REFUSED} & Port Closed & \textbf{NO} & \texttt{Task.status} $\to$ \texttt{IN\_PROGRESS} & Linear Backoff \\
\texttt{MSH-10 MISMATCH} & Correlation Breach & \textbf{NO} & \texttt{Task.status} $\to$ \texttt{FAILED} & Operator Review \\
\bottomrule
\end{tabularx}
\caption{Downstream HL7 Acknowledgement Code Decision Matrix}
\label{tab:egress_ack_decision}
\end{table}

\section{State, Task, and Provenance Lifecycle Management}
\label{sec:egress_lifecycle}

\subsection{Outbound State Lifecycle Synchronization}

Upon receiving the \texttt{OutboundMllpResponse}, the \texttt{OutboundStateLifecycleManager} updates the associated canonical FHIR resources in \texttt{Mneme} and \texttt{Mnemosyne}, constructs an immutable \texttt{Provenance} record, and emits an audit event to Petasos:

\begin{lstlisting}[language=Java, caption={OutboundStateLifecycleManager Implementation}]
@ApplicationScoped
public class OutboundStateLifecycleManager {

    @Inject
    private TaskService taskService;

    @Inject
    private CommunicationService communicationService;

    @Inject
    private ProvenanceService provenanceService;

    @Inject
    private OutboundProvenanceResourceBuilder provenanceBuilder;

    @Inject
    private TaskEventProducerService taskEventProducerService;

    public void handleOutboundResponse(OutboundMllpRequest request, OutboundMllpResponse response) {
        // 1. Update Synthetic Task (Pragma) Status
        if (StringUtils.isNotBlank(request.getTaskId())) {
            Task task = taskService.getTask(request.getTaskId());
            if (task != null) {
                if (response.isSuccessful()) {
                    task.setStatus(Task.TaskStatus.COMPLETED);
                    task.setBusinessStatus(new CodeableConcept().setText("DELIVERED_AA"));
                } else if ("AE".equals(response.getAckCode())) {
                    task.setStatus(Task.TaskStatus.INPROGRESS);
                    task.setBusinessStatus(new CodeableConcept().setText("RETRY_PENDING_AE"));
                } else {
                    task.setStatus(Task.TaskStatus.FAILED);
                    task.setBusinessStatus(new CodeableConcept().setText("REJECTED_AR"));
                }
                task.setLastModified(new Date());
                taskService.update(task);
            }
        }

        // 2. Build and Persist Immutable Provenance Record
        Provenance provenance = provenanceBuilder.buildProvenance(request, response);
        provenanceService.save(provenance);

        // 3. Publish Egress Completion Event to Petasos
        ErgonEvent completionEvent = new ErgonEvent();
        completionEvent.setEventId(UUID.randomUUID().toString());
        completionEvent.setTaskId(request.getTaskId());
        completionEvent.setDestinationId(response.getDestinationId());
        completionEvent.setStatus(response.isSuccessful() ? "SUCCESS" : "FAILURE");
        completionEvent.setDurationMs(response.getDurationMs());

        taskEventProducerService.publishOutboundCompletionEvent("petasos.queue.egress.completed", completionEvent);
    }
}
\end{lstlisting}

\subsection{FHIR Provenance Resource Construction}

The \texttt{OutboundProvenanceResourceBuilder} captures complete transmission audit trails:

\begin{lstlisting}[language=Java, caption={OutboundProvenanceResourceBuilder Blueprint}]
public class OutboundProvenanceResourceBuilder {

    public Provenance buildProvenance(OutboundMllpRequest request, OutboundMllpResponse response) {
        Provenance prov = new Provenance();
        prov.setId(UUID.randomUUID().toString());
        prov.setRecorded(new Date());

        // Target References (Task & Communication)
        if (StringUtils.isNotBlank(request.getTaskId())) {
            prov.addTarget(new Reference("Task/" + request.getTaskId()));
        }
        if (StringUtils.isNotBlank(request.getCommunicationId())) {
            prov.addTarget(new Reference("Communication/" + request.getCommunicationId()));
        }

        // Activity: Data Transmission
        CodeableConcept activity = new CodeableConcept();
        activity.addCoding().setSystem("http://terminology.hl7.org/CodeSystem/v3-DataOperation")
                            .setCode("TRANSMIT")
                            .setDisplay("MLLP Outbound Transmission");
        prov.setActivity(activity);

        // Agent: Transmitter Gateway
        Provenance.ProvenanceAgentComponent agent = new Provenance.ProvenanceAgentComponent();
        agent.getType().addCoding().setSystem("http://terminology.hl7.org/CodeSystem/provenance-participant-type")
                                   .setCode("transmitter")
                                   .setDisplay("Harmonia MLLP Egress Gateway");
        agent.setWho(new Reference().setDisplay("pylai-mllp-out-" + response.getDestinationId()));
        prov.addAgent(agent);

        return prov;
    }
}
\end{lstlisting}

% =============================================================================
% PART III: TESTING, VERIFICATION & CONTAINERIZATION
% =============================================================================
\section{Testing, Verification Harnesses \& Containerization}
\label{sec:mllp_testing_containerization}

To guarantee high availability and fault resilience prior to hospital production rollout, new MLLP ingress and egress services must be thoroughly validated across four testing tiers.

\subsection{Tier 1: Unit Testing with HAPI HL7 v2 Parsers}

Unit tests validate segment parsing, field extraction, and boundary conditions without network dependencies:

\begin{lstlisting}[language=Java, caption={Unit Test Suite for ADT Message Extraction}]
public class AdtMessageExtractorTest {

    private final PipeParser parser = new PipeParser();

    @Test
    public void testSuccessfulAdtExtraction() throws Exception {
        String hl7 = "MSH|^~\\&|PAS|HOSPITAL|HARMONIA|HIE|20260916080000||ADT^A01|MSG-100421|P|2.4\r"
                   + "PID|1||MRN-90210^^^HOSP||DOE^JANE^^^^||19850412|F\r"
                   + "PV1|1|I|WARD-3A^RM-12^BED-02||||1234^SMITH^JOHN^^DR|||||||||||VIS-88001";

        Message msg = parser.parse(hl7);
        AdtMessageExtractor extractor = new AdtMessageExtractor(msg);

        assertEquals("MSG-100421", extractor.extractMessageControlId());
        assertEquals("A01", extractor.extractTriggerEvent());
        assertEquals("MRN-90210", extractor.extractPatientId());
        assertEquals("DOE", extractor.extractFamilyName());
        assertEquals("JANE", extractor.extractGivenName());
        assertEquals("VIS-88001", extractor.extractVisitNumber());
    }
}
\end{lstlisting}

\subsection{Tier 2: Camel Route Harness Testing}

Camel routes are validated using \texttt{CamelTestSupport} and \texttt{MockEndpoint} assertions:

\begin{lstlisting}[language=Java, caption={Camel Route Mock Test Harness}]
public class OutboundMllpRouteTest extends CamelTestSupport {

    @EndpointInject("mock:mllp-target")
    private MockEndpoint mockTarget;

    @Test
    public void testSuccessfulOutboundDispatch() throws Exception {
        mockTarget.expectedMessageCount(1);
        mockTarget.whenAnyExchangeReceived(exchange -> {
            exchange.getMessage().setBody("MSH|^~\\&|||||20260916080001||ACK^A01|ACK-01|P|2.4\rMSA|AA|MSG-100421");
        });

        OutboundMllpRequest req = new OutboundMllpRequest();
        req.setMessageControlId("MSG-100421");
        req.setDestinationId("EMR");
        req.setRawHl7Payload("MSH|^~\\&|HARMONIA||EMR||20260916080000||ADT^A01|MSG-100421|P|2.4\rPID|1||MRN-90210");

        OutboundMllpResponse resp = template.requestBody("direct:mllp-outbound-send", req, OutboundMllpResponse.class);

        assertTrue(resp.isSuccessful());
        assertEquals("AA", resp.getAckCode());
        mockTarget.assertIsSatisfied();
    }
}
\end{lstlisting}

\subsection{Tier 3: Embedded Mock MLLP Server Integration Testing}

Integration tests launch embedded TCP servers (via Paradeigma's \texttt{MllpServer}) to simulate real network conditions, latency, and connection resets:

\begin{lstlisting}[language=Java, caption={Embedded MLLP Integration Test with Fault Simulation}]
@Test
public void testTimeoutAndApplicationRejectHandling() throws Exception {
    // Launch embedded mock MLLP receiver returning AR Reject
    MllpServer mockHospitalServer = new MllpServer(2204, (req) -> {
        return "MSH|^~\\&|||||20260916080002||ACK^O01|ACK-99|P|2.4\rMSA|AR|ORD-5501|Patient MRN Not Found";
    });
    mockHospitalServer.start();

    try {
        OutboundMllpRequest req = new OutboundMllpRequest();
        req.setMessageControlId("ORD-5501");
        req.setDestinationId("LMS");
        req.setTargetPort(2204);
        req.setRawHl7Payload("MSH|^~\\&|HARMONIA||LMS||...|ORM^O01|ORD-5501|P|2.4\r...");

        OutboundMllpResponse resp = template.requestBody("direct:mllp-outbound-send", req, OutboundMllpResponse.class);

        assertFalse(resp.isSuccessful());
        assertEquals("AR", resp.getAckCode());
        assertTrue(resp.getErrorMessage().contains("Patient MRN Not Found"));
    } finally {
        mockHospitalServer.stop();
    }
}
\end{lstlisting}

\subsection{Tier 4: Multi-Instance Docker Compose Deployment Topology}

Production deployments isolate egress gateways into independent containers with resource quotas and dedicated queues:

\begin{lstlisting}[language=bash, caption={Docker Compose Multi-Instance Egress Architecture}]
version: '3.8'

services:
  # Ingress Gateway (Inbound MLLP Listeners)
  pylai-mllp-in:
    image: harmonia/pylai-mllp-in:1.0.0
    ports:
      - "2101:2101" # ADT Inbound
      - "2102:2102" # Lab ORU Inbound
      - "2103:2103" # Rad ORU Inbound
      - "2104:2104" # ORM Inbound
    environment:
      - MLLP_BIND_HOST=0.0.0.0
      - ARTEMIS_BROKER_URL=tcp://petasos-broker:61616
      - INFINISPAN_HOTROD_SERVERS=mneme-cluster:11222

  # Egress Gateway Instance 1: Dedicated EMR Gateway
  pylai-mllp-out-emr:
    image: harmonia/pylai-mllp-out:1.0.0
    environment:
      - MLLP_OUTBOUND_DESTINATION_ID=EMR
      - MLLP_OUTBOUND_HOST=emr.hospital.internal
      - MLLP_OUTBOUND_PORT=2201
      - MLLP_OUTBOUND_QUEUE_NAME=petasos.queue.egress.emr.adt
      - MLLP_OUTBOUND_READ_TIMEOUT_MS=8000
      - ARTEMIS_BROKER_URL=tcp://petasos-broker:61616

  # Egress Gateway Instance 2: Dedicated LMS Gateway
  pylai-mllp-out-lms:
    image: harmonia/pylai-mllp-out:1.0.0
    environment:
      - MLLP_OUTBOUND_DESTINATION_ID=LMS
      - MLLP_OUTBOUND_HOST=lms.lab.hospital.internal
      - MLLP_OUTBOUND_PORT=2204
      - MLLP_OUTBOUND_QUEUE_NAME=petasos.queue.egress.lms.orm
      - MLLP_OUTBOUND_READ_TIMEOUT_MS=15000
      - ARTEMIS_BROKER_URL=tcp://petasos-broker:61616

  # Egress Gateway Instance 3: Dedicated RIS-PACS Gateway
  pylai-mllp-out-rispac:
    image: harmonia/pylai-mllp-out:1.0.0
    environment:
      - MLLP_OUTBOUND_DESTINATION_ID=RISPAC
      - MLLP_OUTBOUND_HOST=rispac.radiology.hospital.internal
      - MLLP_OUTBOUND_PORT=2205
      - MLLP_OUTBOUND_QUEUE_NAME=petasos.queue.egress.rispac.orm
      - MLLP_OUTBOUND_READ_TIMEOUT_MS=20000
      - ARTEMIS_BROKER_URL=tcp://petasos-broker:61616
\end{lstlisting}

\newpage
